\indent ಈ ಉದಾಹರಣೆಗಳಲ್ಲಿ ಏಕವಚನ ಶಬ್ದಗಳಿಗೆ “ಗಳು” ಎಂಬ ಆಗಮ ಸೇರಿ ಬಹುವಚನ ಶಬ್ದಗಳಾಗಿರುವುದು ಕಂಡುಬರುತ್ತದೆ. ಆಳು, ಹಸು, ರೂಪಾಯಿ, ಗಿಣಿ—ಇವು ನಪುಂಸಕ ಲಿಂಗ ಶಬ್ದಗಳು. ಆದುದರಿಂದ ನಪುಂಸಕ ಲಿಂಗ ಶಬ್ದಗಳಿಗೆ ಬಹುವಚನದಲ್ಲಿ “ಗಳು” ಎಂಬ ಆಗಮ ಸೇರುತ್ತದೆಂದು ತಿಳಿಯಬಹುದು.\par

\indent ವಿಶೇಷ ಸಂದರ್ಭ : ಅಕಾರಾಂತವಲ್ಲದ ಪುಲ್ಲಿಂಗ ಹಾಗೂ ಸ್ತ್ರೀಲಿಂಗ ಶಬ್ದಗಳಿಗೆ “ಗಳು” ಎಂಬ ಆಗಮ ಸೇರಿ ಬಹುವಚನ ರೂಪಗಳಾಗುವುದು.\par

\vspace{3pt plus 2pt}\noindent\hspace*{1.2em}\kbvlabel{ಉದಾ}\par\vspace{1pt}

\begin{kbvderiv}
\vspace{4pt plus 3pt minus 1pt}
\begin{center}\begin{adjustbox}{max width=\textwidth}
{\renewcommand{\arraystretch}{1.22}\setlength{\tabcolsep}{4pt}
\begin{tabular}{l c l}
ಗಂಡು + ಗಳು & {\color{kbvgrey}$=$} & \kbvresult{ಗಂಡುಗಳು} \\
ಹೆಣ್ಣು + ಗಳು & {\color{kbvgrey}$=$} & \kbvresult{ಹೆಣ್ಣುಗಳು} \\
\end{tabular}}
\end{adjustbox}\end{center}
\vspace{4pt plus 3pt minus 1pt}
\end{kbvderiv}

\begin{kbvrule}
\noindent\kbvbullet\hspace{0.5em}ಏಕವಚನದಲ್ಲಿರುವ ಪುಲ್ಲಿಂಗ, ಸ್ತ್ರೀಲಿಂಗ ಶಬ್ದಗಳಿಗೆ ಸಾಮಾನ್ಯ ವಾಗಿ “ಅರು” ಎಂಬ ಆಗಮ ಸೇರಲು ಬಹುವಚನ ರೂಪವುಂಟಾಗುವುದು.
\end{kbvrule}

\vspace{3pt plus 2pt}\noindent\hspace*{1.2em}\kbvlabel{ಉದಾ}\par\vspace{1pt}

\begin{kbvderiv}
\vspace{4pt plus 3pt minus 1pt}
\begin{center}\begin{adjustbox}{max width=\textwidth}
{\renewcommand{\arraystretch}{1.22}\setlength{\tabcolsep}{4pt}
\begin{tabular}{l c l}
ಅರಸ + ಅರು & {\color{kbvgrey}$=$} & \kbvresult{ಅರಸರು}{\color{kbvgrey} (ಬ.ವ.)} \\
ಹುಡುಗ + ಅರು & {\color{kbvgrey}$=$} & \kbvresult{ಹುಡುಗರು} \\
ಮುದುಕಿ + ಅರು & {\color{kbvgrey}$=$} & \kbvresult{ಮುದುಕಿಯರು} \\
ತಂಗಿ + ಅರು & {\color{kbvgrey}$=$} & \kbvresult{ತಂಗಿಯರು} \\
\end{tabular}}
\end{adjustbox}\end{center}
\vspace{4pt plus 3pt minus 1pt}
\end{kbvderiv}

\begin{kbvrule}
\noindent\kbvbullet\hspace{0.5em}ಏಕವಚನದಲ್ಲಿರುವ ಕೆಲವು ಪುಲ್ಲಿಂಗ, ಸ್ತ್ರೀಲಿಂಗ ಅಕಾರಾಂತ ಬಂಧುವಾಚಕ ಶಬ್ದಗಳಿಗೆ “ಅಂದಿರು” ಎಂಬ ಆಗಮ ಸೇರಲು ಬಹುವಚನ ರೂಪವುಂಟಾಗುವುದು.
\end{kbvrule}

\vspace{3pt plus 2pt}\noindent\hspace*{1.2em}\kbvlabel{ಉದಾ}\par\vspace{1pt}

\begin{kbvderiv}
\vspace{4pt plus 3pt minus 1pt}
\begin{center}\begin{adjustbox}{max width=\textwidth}
{\renewcommand{\arraystretch}{1.22}\setlength{\tabcolsep}{4pt}
\begin{tabular}{l c l}
ಅಣ್ಣ + ಅಂದಿರು & {\color{kbvgrey}$=$} & \kbvresult{ಅಣ್ಣಂದಿರು}{\color{kbvgrey} (ಬ.ವ.)} \\
ಅಳಿಯ + ಅಂದಿರು & {\color{kbvgrey}$=$} & \kbvresult{ಅಳಿಯಂದಿರು} \\
ಅಕ್ಕ + ಅಂದಿರು & {\color{kbvgrey}$=$} & \kbvresult{ಅಕ್ಕಂದಿರು} \\
ತಾಯಿ + ಅಂದಿರು & {\color{kbvgrey}$=$} & \kbvresult{ತಾಯಂದಿರು} \\
\end{tabular}}
\end{adjustbox}\end{center}
\vspace{4pt plus 3pt minus 1pt}
\end{kbvderiv}

\begin{kbvrule}
\noindent\kbvbullet\hspace{0.5em}ಕೆಲವು ಏಕವಚನ ಶಬ್ದಗಳಿಗೆ ಕಳು, ವು ಮೊದಲಾದ ಆಗಮಗಳು ಸೇರಿ ಬಹುವಚನ ರೂಪಗಳಾಗುವುದು.
\end{kbvrule}
